% 02_modelado_estado.tex - Capítulo 2: Modelado del Problema y Representación de Estados
\chapter{Modelado del Problema y Representación de Estados}
\label{ch:modelado}

Este capítulo describe el modelado del problema como espacio de estados (Unidad 2 del curso) y
la representación del estado compuesto que soporta tanto las búsquedas de un agente (F2, F3)
como la búsqueda adversarial (F4). Cumple con \rf{10}--\rf{13} y el caso 6.

\section{Principios de Diseño}
\label{sec:principios-diseno}

El modelado sigue tres principios fundamentales:

\begin{enumerate}
\item \textbf{Inmutabilidad:} \code{GameState} y \code{Board} son \code{@dataclass(frozen=True)}.
    La inmutabilidad garantiza que los estados puedan usarse como claves en \code{dict}/\code{set}
    (visited, frontier, parents) sin riesgo de mutación accidental.
\item \textbf{Separación estática/dinámica:} \code{Board} contiene todo lo que \textbf{nunca
    cambia} (grid, catálogo de terrenos, bases, bando de cada unidad, límite de turnos).
    \code{GameState} contiene solo lo que \textbf{cambia} (posiciones, carrier, resource, turn,
    turn\_number).
\item \textbf{Estado compuesto completo (RF-10, RF-11):} Un estado no es solo una posición;
    dos estados con la misma posición física son distintos si difieren en \code{carrier},
    \code{resource}, \code{turn}, \code{turn\_number} o las posiciones de otras unidades.
    El     \textbf{caso 6} demuestra esto explícitamente (ver \tableref{tab:caso6}).
\end{enumerate}

\section{Partes Estáticas: \code{Board}}
\label{sec:board}

\code{Board} es la parte estática del escenario; se construye una sola vez al cargar el JSON y
nunca muta. Contiene:

\begin{itemize}
\item \code{rows}, \code{cols}: dimensiones del grid
\item \code{grid}: tupla de tuplas con nombres de terreno por celda
\item \code{terrains}: \code{Mapping[str, Terrain]} -- catálogo \code{nombre $\to$ Terrain(cost, passable)}
\item \code{bases}: \code{Mapping[Side, Position]} -- base de cada bando
\item \code{unit\_sides}: \code{Mapping[str, Side]} -- bando de cada unidad (clave: id)
\item \code{turn\_limit}: límite de turnos para empate (default 100, configurable en JSON)
\end{itemize}

Métodos clave:
\begin{minted}{python}
def in_bounds(self, pos: Position) -> bool
def terrain_at(self, pos: Position) -> Terrain
def is_passable(self, pos: Position) -> bool
def cost_at(self, pos: Position) -> float
\end{minted}

\section{Parte Dinámica: \code{GameState}}
\label{sec:gamestate}

\code{GameState} es el estado compuesto inmutable (frozen dataclass). Contiene:

\begin{minted}{python}
@dataclass(frozen=True)
class GameState:
    positions: tuple[tuple[str, Position], ...]  # ordenado por unit_id
    carrier: str | None          # id de la unidad que lleva el recurso
    resource: Position | None    # celda del recurso si nadie lo lleva
    turn: Side                   # Side.A o Side.B
    turn_number: int = 0         # para límite de turnos / empate
\end{minted}

Métodos auxiliares:
\begin{minted}{python}
def position_of(self, unit_id: str) -> Position
def unit_at(self, pos: Position) -> str | None
\end{minted}

\section{Posición y Vecindad}
\label{sec:position}

\code{Position} es una tupla inmutable \code{(row, col)} con métodos:

\begin{minted}{python}
def neighbors(self) -> tuple[Position, ...]:  # 4 dirs: arriba, abajo, izq, der
def is_adjacent(self, other: Position) -> bool  # distancia Manhattan == 1
\end{minted}

La vecindad es de \textbf{4 direcciones} (arriba, abajo, izquierda, derecha) por decisión
del equipo (decisión 2026-10-07). Con 4 direcciones, la heurística Manhattan es admisible.

\section{Caso 6: Estado Compuesto Demostrado}
\label{sec:caso6}

El \textbf{caso 6} demuestra \rf{11}: ``dos estados con una unidad en la misma casilla pueden
ser distintos''. En el mapa del caso 7 (5$\times$7, coste 1), la unidad A1 en (1,0) junto a su
base (0,0):

\begin{table}[H]
\centering
\caption{Caso 6: mismo posición, estados distintos}
\begin{tabular}{lllll}
\toprule
\textbf{Estado} & \textbf{Pos A1} & \textbf{Carrier} & \textbf{Turn} & \textbf{Efecto de Move A1 $\to$ (0,0)} \\
\midrule
Libre (inicial) & (1,0) & nadie & A & entra a la base, no hay ganador \\
Portando & (1,0) & A1 & A & entra a la base $\to$ \textbf{gana A} \\
Turno rival & (1,0) & nadie & B & A1 no tiene acciones legales \\
\bottomrule
\end{tabular}
\end{table}
\label{tab:caso6}

Los tres estados tienen \code{positions} idénticas pero son objetos distintos
(\code{len(set(states)) == 3}). Un diccionario indexado por estado mantiene tres entradas;
indexado por posición solo mantendría una. Test: \texttt{tests/test\_case\_6.py}.

\section{Identificación Unívoca (RF-12)}
\label{sec:identificacion-univoca}

La Tabla~\ref{tab:identificacion} resume dónde se encuentra cada elemento requerido por \rf{12}.

\begin{table}[H]
\centering
\caption{Identificación unívoca de elementos del problema (RF-12)}
\begin{tabular}{ll}
\toprule
\textbf{Elemento} & \textbf{Dónde se define} \\
\midrule
Estado inicial & \code{Scenario.initial\_state} (construido por \code{scenario/loader.py}) \\
Acciones válidas & \code{domain.rules.legal\_actions(board, state)} \\
Estados sucesores & \code{domain.rules.successors(board, state)} $\to$ \code{(action, next\_state, cost)} \\
Condición de meta & Búsqueda: \code{problem.is\_goal(state)}; Juego: \code{winner(board, state)} \\
Costo de acción & \code{domain.rules.apply(board, state, action)} $\to$ \code{(next\_state, cost)} \\
\bottomrule
\end{tabular}
\label{tab:identificacion}
\end{table}

\section{Justificación de Suficiencia (RF-13)}
\label{sec:justificacion-suficiencia}

El estado contiene \textbf{exactamente} la información necesaria para decidir:

\begin{enumerate}
\item \textbf{Dónde está cada unidad} (\code{positions}) $\to$ legalidad de movimientos, ocupación
\item \textbf{De qué bando es cada unidad} (\code{unit\_sides} en \code{Board}) $\to$ turno, legalidad
\item \textbf{Quién lleva el recurso} (\code{carrier}) $\to$ pickup, delivery, interceptación
\item \textbf{Dónde está el recurso si nadie lo lleva} (\code{resource}) $\to$ pickup
\item \textbf{De quién es el turno} (\code{turn}) $\to$ qué unidades pueden moverse
\item \textbf{Número de turno} (\code{turn\_number}) $\to$ empate por límite
\end{enumerate}

Nada más influye en las reglas (no hay vida, munición, información oculta). Por tanto, dos
estados con iguales valores en estos seis campos tienen idéntico conjunto de acciones legales,
mismos sucesores, mismos costos y mismo resultado final. \rf{13} satisfecho.